\chapter{Introduction}

Specification sheets quote cache latencies in cycles, memory bandwidth in
gigabytes per second, and peak floating point throughput in GFLOPS. These
numbers are ceilings under idealized assumptions. What a program actually sees
depends on the access pattern, the compiler, the memory hierarchy under load,
and the clock the core really runs at. This project measures those numbers
directly on one machine, an Intel Core i7-14700K, and assembles them into a
model that predicts something concrete: the cache block sizes a blocked matrix
multiply should use.

The work has five measured pillars. First, a dependent load pointer chase
recovers the load use latency of each cache level, from a 4 KB working set that
lives in L1 up to a 512 MB working set that is firmly in DRAM, and a change
point detector locates the plateau boundaries and cross checks them against the
cache sizes the operating system reports. Second, a STREAM style benchmark in
three variants measures sustainable bandwidth and a thread scaling curve.
Third, a ten chain fused multiply add loop measures peak single precision
throughput and compares it to the ceiling computed from the measured clock.
Fourth, these ceilings become a roofline, and a BLIS style analytical model
predicts matrix multiply tile sizes that a blocked GEMM sweep then tests.
Fifth, the compute kernels are ported to AArch64 NEON and checked for
correctness under emulation.

The guiding principle is measurement honesty. The target runs inside WSL2,
which exposes no hardware performance counters, so every number here comes from
timing and a known transfer size, the classical approach of
McCalpin~\cite{mccalpin1995}. Where the platform makes a number hard to get
right, this report states the difficulty and the workaround rather than quietly
publishing a plausible figure. Two such difficulties, described in the
methodology, turned out to matter a great deal: the reported clock does not
track the boost frequency, and the hybrid core layout is invisible to the
guest.
